\documentclass{article}

\usepackage[serbian]{babel}
\usepackage[T1]{fontenc}
\usepackage[utf8]{inputenc}

\usepackage{amsmath,amssymb}
\usepackage{graphicx}
\usepackage{caption}
\usepackage{subcaption}
\usepackage{geometry}
\usepackage{hyperref}
\usepackage{amsthm}
\newtheorem{lemma}{Lema}
\newtheorem{theorem}{Teorema}



\hypersetup{
    colorlinks=true,
    linkcolor=blue,
    citecolor=blue,
    urlcolor=blue
}

\title{Minimizacija Menhetn rastojanja}
\author{Staša Đorđević 1007/2025\\
\small Matematički fakultet Univerziteta u Beogradu}
\date{}

\begin{document}

\maketitle


\section{Uvod}

Neka je dat skup tačaka $P = \{p_1, p_2, \dots, p_n\}$ u ravni, gde svaka tačka 
$p_i = (x_i, y_i)$ ima celobrojne koordinate. 
Za dve tačke $p_i$ i $p_j$ njihovo \emph{Menhetn rastojanje} definisano je formulom

\[
d(p_i, p_j) = |x_i - x_j| + |y_i - y_j|.
\]

Menhetn rastojanje predstavlja dužinu najkraće putanje
između dve tačke ukoliko je kretanje dozvoljeno samo horizontalno i vertikalno.

U ovom radu razmatramo sledeći problem:  
za dati skup tačaka potrebno je ukloniti tačno jednu tačku tako da maksimalno
Menhetn rastojanje između bilo koje dve preostale tačke bude minimalno.

Formalno, za svaku tačku $p_k \in P$ posmatramo skup
\[
P_k = P \setminus \{p_k\}
\]
i definišemo
\[
D_k = \max_{p_i, p_j \in P_k} d(p_i, p_j).
\]
Zadatak je pronaći
\[
\min_{1 \le k \le n} D_k.
\]

Intuitivno, cilj je ukloniti tačku koja najviše razvlači skup, odnosno
tačku čije prisustvo uzrokuje veliko maksimalno rastojanje između nekog para tačaka.

\subsection*{Primer}

Posmatrajmo skup od četiri tačke:
\[
(0,0),\ (1,1),\ (2,2),\ (10,10).
\]
Tačka $(10,10)$ je značajno udaljena od ostalih i određuje maksimalno
Menhetn rastojanje u skupu.  
Uklanjanjem ove tačke maksimalno rastojanje među preostalim tačkama
značajno se smanjuje, što ilustruje suštinu problema.

\medskip

\section{Naivno rešenje}


Naivni pristup podrazumeva da za svaku tačku računamo Menhetn rastojanje do svih
ostalih tačaka u skupu, čime možemo odrediti maksimalno rastojanje i indekse
tačaka koje ga ostvaruju.

Nakon toga, potrebno je odlučiti koju od te dve tačke ukloniti.
To možemo uraditi tako što za svaku od njih posebno izračunamo maksimalno
Menhetn rastojanje u skupu koji se dobija izbacivanjem druge tačke.
Tačka čijim uklanjanjem dobijamo manju maksimalnu udaljenost predstavlja
bolji izbor za uklanjanje.

Međutim, moguće je da već postoji par tačaka sa većim rastojanjem koji ne
uključuje nijednu od prethodno posmatranih tačaka. U tom slučaju,
izbor između te dve tačke nije od presudnog značaja, jer njihovo uklanjanje
ne utiče na globalni maksimum.

Alternativno, sve izračunate udaljenosti možemo sačuvati u pogodnoj
strukturi podataka koja preslikava vrednosti rastojanja u indekse tačaka
koje ih ostvaruju, pri čemu je struktura uređena po veličini rastojanja.
Iz takve strukture možemo najpre uzeti najveću vrednost i pročitati indekse
tačaka $i$ i $j$ koje je ostvaruju. Zatim posmatramo sledeću najveću
vrednost. Ako ona uključuje neku od tačaka $i$ ili $j$, tu tačku biramo
za uklanjanje. Ukoliko druga najveća vrednost ne uključuje nijednu od njih,
izbor između $i$ i $j$ nije bitan.

Ovakav pristup zahteva proveru svih parova tačaka, pa je njegova vremenska
složenost $O(n^2)$, što nije optimalno za veće ulazne skupove.

\section{Optimalno rešenje}

Posmatrajmo dve tačke $p_i = (x_i, y_i)$ i $p_j = (x_j, y_j)$.
Njihovo Menhetn rastojanje dato je formulom

\[
d(p_i,p_j) = |x_i - x_j| + |y_i - y_j|.
\]

Zavisno od međusobnog položaja tačaka u ravni, apsolutne vrednosti se mogu
razložiti na četiri slučaja:

\[
d(p_i,p_j) =
\begin{cases}
(x_i - x_j) + (y_i - y_j), & x_i \ge x_j,\ y_i \ge y_j, \\
(x_i - x_j) - (y_i - y_j), & x_i \ge x_j,\ y_i < y_j, \\
-(x_i - x_j) + (y_i - y_j), & x_i < x_j,\ y_i \ge y_j, \\
-(x_i - x_j) - (y_i - y_j), & x_i < x_j,\ y_i < y_j.
\end{cases}
\]

Grupisanjem članova koji zavise od indeksa $i$ i $j$ dobijamo:

\[
d(p_i,p_j) =
\begin{cases}
(x_i + y_i) - (x_j + y_j), \\
(x_i - y_i) - (x_j - y_j), \\
-(x_i - y_i) + (x_j - y_j), \\
-(x_i + y_i) + (x_j + y_j).
\end{cases}
\]

Odavde vidimo da se javljaju samo sume i razlike koordinata jedne tačke.
Uvedimo sada sledeće oznake:

\[
\text{sum}_i = x_i + y_i, \qquad
\text{diff}_i = x_i - y_i.
\]

Tada se Menhetn rastojanje može zapisati u kompaktnijem obliku:

\[
d(p_i,p_j) =
\begin{cases}
\text{sum}_i - \text{sum}_j, \\
\text{diff}_i - \text{diff}_j, \\
-(\text{diff}_i - \text{diff}_j) = \text{diff}_j - \text{diff}_i,  \\
-(\text{sum}_i - \text{sum}_j) = \text{sum}_j - \text{sum}_i.
\end{cases}
\]

Iz prethodnog izraza sledi da maksimalno Menhetn rastojanje između bilo koje
dve tačke u skupu može biti ostvareno ili kao maksimalna razlika vrednosti
$\text{sum}_i$, ili kao maksimalna razlika vrednosti $\text{diff}_i$.


Preciznije, važi

\[
\max_{i,j} d(p_i,p_j)
=
\max \left(
\max_i \text{sum}_i - \min_i \text{sum}_i,\ 
\max_i \text{diff}_i - \min_i \text{diff}_i
\right).
\]

Dakle, problem pronalaženja maksimalnog Menhetn rastojanja svodi se na
pronalaženje ekstremnih vrednosti transformisanih koordinata
$\text{sum}_i = x_i + y_i$ i $\text{diff}_i = x_i - y_i$, što omogućava
rešavanje u linearnom vremenu.

Kad smo našli maksimalno rastojanje (neka ono bude između tačaka sa indeksima i i j), treba videti koju od dve tačke treba ukloniti tako da minimizujemo maksimalno rastojanje skupa tačaka. Jasno je da će uklanjanjem tačno jedne od njih maksimalno rastojanje biti manje ili jednako, nikako veće. Treba proveriti sada drugo najveće rastojanje i tačke koje učestvuju u njemu. To možemo uraditi na isti način kao što smo odredili maksimalno rastojanje. Ako jedna od tačaka $p_i$ ili $p_j$ učestvuje u drugom maksimalnom rastojanju, nju treba ukloniti. Inače je svejedno koja će biti uklonjena.

Vremenska složenost ovakvog pristupa je $O(n)$.

\section{Implementacija}

\begin{thebibliography}{9}

\bibitem{rad}
Timothy M. Chan, Nan Hu. \emph{Geometric red–blue set cover for unit squares and related
problems} Journal of Computational Geometry (2014): 380-385.

\end{thebibliography}

\end{document}